\documentclass[11pt,a4paper]{article}
\usepackage[T1]{fontenc}
\usepackage[utf8]{inputenc}
\usepackage{lmodern,microtype}
\usepackage[a4paper,left=25mm,right=25mm,top=24mm,bottom=26mm,headheight=15pt]{geometry}
\usepackage{amsmath,amssymb,amsthm,mathtools}
\usepackage{booktabs,array,longtable,enumitem}
\usepackage[dvipsnames]{xcolor}
\usepackage{fancyhdr,needspace}
\usepackage[unicode,colorlinks=true,linkcolor=MidnightBlue,citecolor=MidnightBlue,urlcolor=MidnightBlue,breaklinks=true]{hyperref}
\usepackage{bookmark,xurl}
\pagestyle{fancy}
\fancyhf{}
\fancyhead[L]{\small Exact type beyond log-convexity}
\fancyhead[R]{\small Formal nonlinear calculus}
\fancyfoot[C]{\thepage}
\renewcommand{\headrulewidth}{0.3pt}
\setlength{\parskip}{3pt}
\setlength{\emergencystretch}{2em}
\setlist{leftmargin=1.7em,itemsep=3pt,topsep=5pt}
\numberwithin{equation}{section}
\newtheorem{theorem}{Theorem}[section]
\newtheorem{lemma}[theorem]{Lemma}
\newtheorem{proposition}[theorem]{Proposition}
\newtheorem{corollary}[theorem]{Corollary}
\theoremstyle{definition}
\newtheorem{definition}[theorem]{Definition}
\newtheorem{example}[theorem]{Example}
\theoremstyle{remark}
\newtheorem{remark}[theorem]{Remark}
% ed. (2026-09-29): unnumbered environment for ProveIt editorial notes; remarks share the theorem counter, so the notes must not be numbered.
\newtheorem*{ednote}{Editorial note (ProveIt, 2026-09-29)}
\newcommand{\R}{\mathbb R}
\newcommand{\C}{\mathbb C}
\newcommand{\N}{\mathbb N}
\newcommand{\Q}{\mathbb Q}
\newcommand{\id}{\operatorname{id}}
\newcommand{\Sym}{\operatorname{Sym}}
\newcommand{\Res}{\operatorname{Res}}
\newcommand{\Type}{T_N}
\newcommand{\inv}[1]{#1^{\langle-1\rangle}}
\newcommand{\coeff}[2]{[#1^{#2}]}
\newcommand{\norm}[1]{\left\lVert#1\right\rVert}
\newcommand{\ee}{\mathrm e}
\newcommand{\dd}{\mathrm d}
\newcommand{\lp}{\log^{+}}
\newcommand{\ent}{\mathfrak h}
\newcommand{\snapshot}{\nolinkurl{3d5973524506411392a911470b5ddc35521568ea}}
\hypersetup{pdftitle={Exact Type Beyond Log-Convexity},pdfauthor={AI-assisted research prepared for Vladimir Reshetnikov},pdfsubject={Complete formal reversion criterion, supermultiplicative envelopes, optimal distortion, divergent composition, Banach-space formal maps}}
\begin{document}
\hypersetup{pageanchor=false}
\begin{titlepage}
\centering
\vspace*{11mm}
{\Large\scshape Transseries and nonlinear inversion\par}
\vspace{12mm}
{\Huge\bfseries Exact Type Beyond\\[3mm] Log-Convexity\par}
\vspace{9mm}
{\Large A complete reversion criterion, sharp distortion,\\[2mm]
and dimension-free nonlinear calculus\par}
\vspace{12mm}
{\large Prepared for Vladimir Reshetnikov\par}
\vspace{3mm}
29 September 2026
\vspace{11mm}
\begin{minipage}{0.91\textwidth}
\begin{abstract}
For an arbitrary positive coefficient weight $N=(N_n)$, we characterize
exact preservation of exponential coefficient type under every formal
compositional inverse. The criterion is superexponential root growth
\emph{together} with subexponential agreement with the least supermultiplicative
majorant
\[
 S_n=\max_{j_1+\cdots+j_k=n}\prod_{i=1}^{k}N_{j_i}.
\]
The central estimate is stronger: whenever $N_n^{1/n}\to\infty$, the
positive extremal inverse has coefficient $S_n\exp(o(n))$, for every
fixed positive input amplitude. Consequently
$\limsup(S_n/N_n)^{1/n}$ is the exact optimal multiplicative distortion
of type, not merely an upper bound.

For the zero-loss weights we prove a sharp composition inequality for two
divergent arguments, determine all possible cancellations at equal type,
and extend inversion and composition to formal Banach-space maps without
a dimension-dependent loss. Explicit oscillating weights show that
log-convex approximation is not necessary; other examples realize every
finite distortion and even send type zero to infinite type under inversion.
The proofs are self-contained after elementary formal-series algebra.
An exact finite verification program, a source audit, and a further
research program accompany the article.
\end{abstract}
\end{minipage}
\vfill
\begin{minipage}{0.91\textwidth}\small
\textbf{Status and scope.} Conventional mathematical proofs addressing
explicit questions in ProveIt's \emph{Sharp Weighted Type Under Formal
Reversion}. Not Lean-verified or independently peer-reviewed. The
literature comparison is targeted; global publication priority is not
asserted. These are formal coefficient theorems, not analytic summability
theorems.
\end{minipage}
\end{titlepage}
\hypersetup{pageanchor=true}
\setcounter{tocdepth}{1}
\tableofcontents
\vspace{7mm}
\noindent\textbf{Reading guide.}
The central argument is Sections~\ref{sec:envelope}--\ref{sec:classification}.
Section~\ref{sec:composition} gives the complete scalar type calculus;
Section~\ref{sec:banach} specifies the norms needed for its vector-valued
extension. The examples in Section~\ref{sec:examples} separate the exact
criterion from several tempting regularity assumptions. Verification,
formalization, and genuinely remaining questions are treated separately.
\clearpage

\section{Research target and relation to prior work}
\label{sec:context}

\subsection{Three questions from the repository}

The ProveIt article \emph{Sharp Weighted Type Under Formal Reversion}
proves exact preservation of a weighted coefficient limsup for positive
log-convex weights with superexponential roots. Its final research section
asks for a necessary-and-sufficient criterion for arbitrary positive
weights, a sharp rule when both composition arguments are divergent, and
a multivariate or operator-valued extension \cite{repo-weight}. Those are
the specific targets of the present article. We do not revisit the
repository's quadratic-feedback coefficient asymptotics, inverse-digamma
optimal truncation, or Hahn resonance calculations.

The arbitrary-weight question is more than a technical relaxation.
Nonlinear inversion can multiply coefficients whose degrees add to the
requested excess degree. A weight may be large at the input degrees and
have a very deep dip at the output degree. A root-growth condition alone
cannot exclude this failure. Conversely, oscillation need not be harmful:
the products may still fit the target weight with no exponential loss.
The correct condition concerns all finite additive partitions, not
convexity of a sequence of logarithms.

We answer the first question completely for scalar formal power series.
We also give a complete rule based only on the two scalar types: unequal
types force equality in a maximum bound, while every smaller output type
can occur at equal input type. For formal maps tangent to the identity on
a Banach space, the same exact statements hold with multilinear operator
norms. For a general invertible derivative we prove sharp two-sided norm
bounds and explain why the derivative and one scalar input type cannot
determine a single exact output type.

\subsection{What is classical, and what is established here}

The Lagrange inversion formula and its finite coefficient expansions are
classical; Gessel's survey is an appropriate reference \cite{gessel}.
We include the short formal argument needed here and make no novelty
claim about formal invertibility, the coefficient formula, or scalar
majorant methods.

Rainer and Schindl characterize composition in ultradifferentiable
classes \cite{rs-composition}, and characterize equivalences among
composition, inversion, and differential-equation stability under stated
regularity hypotheses \cite{rs-stability}. Their sequence formulation
uses weak log-convexity, derivation conditions, and a Fa\`a di Bruno bound
allowing a factor $C^n$. Their discussion also credits earlier
Banach-space inverse-function results. Thus neither nonlinear class
stability nor Banach-space inversion is a new concept here.

The present question has a different precision and domain. We work with
arbitrary \emph{formal coefficient weights}, retain the exact exponential
limsup constant, and allow neither an unspecified exponential loss nor
an implicit convexity or derivation hypothesis. We prove that one
explicit supermultiplicative envelope measures the entire worst-case
loss. The upper estimate does not assume that the envelope is close to
the original weight. This allows a classification of failure as well as
success, and makes the subsequent composition and vector-valued results
consequences of a single scalar estimate.

The proposed contributions are the envelope asymptotic, the complete
arbitrary-weight criterion and optimal distortion theorem, the resulting
exact scalar type calculus, and the explicit separating examples.
Their proofs are supplied below. A targeted comparison is not an
exhaustive priority search, and the assertions of this article should be
independently checked before being treated as established literature.

\subsection{Snapshot and dependency boundary}

The repository checkpoint used for provenance is
\begin{quote}\small\snapshot.\end{quote}
The principal source is
\begin{quote}\small\ttfamily
Analysis/Transseries/docs/series-and-transseries/\\
Sharp\_Weighted\_Type\_Formal\_Reversion/article.tex.
\end{quote}
Its Git blob is
\begin{quote}\small\nolinkurl{c2d9978806e0861f1bab7108b72ad92773e6a05c}.\end{quote}
In particular, the pinned source's research questions on arbitrary
weights, divergent composition, and multivariate inversion were read
\cite{repo-weight}. Selected neighboring package descriptions were
inspected to avoid obvious overlap. The whole repository, the canonical
large transseries volume, and newly arrived binary research archives were
not audited claim by claim. No upstream file was modified.

No mathematical theorem from an unverified repository article is needed
as a premise of the proofs below. The log-convex result is recovered
independently in Proposition~\ref{prop:logconvex}. No proof assistant was
run; the attached computations are finite checks, not substitutes for
the quantified arguments.
% ed. (2026-09-29): exact location of the three questions at the pinned revision; other questions of that article re-posed in Section 12
\begin{ednote}
At the pinned revision of \cite{repo-weight} the three questions are the
subsections ``Weights beyond log-convexity'', ``Two divergent composition
arguments'' and ``Multivariate and operator-valued inversion'', lines
1485--1493, 1495--1503 and 1505--1513 of
\path{Sharp_Weighted_Type_Formal_Reversion/article.tex};
\path{notes/provenance.json} records the wider range 1457--1547. (Editorial
notes added to that article after this package was filed move these line
numbers; the titles identify the questions.) In Section~\ref{sec:questions},
``Optimal subexponential overhead'', ``Hahn supports and accumulating
actions'', ``Realization and summation compatible with the envelope'' and
the formalization question refine that article's own questions ``Optimal
subexponential cost of reversion'' (pinned lines 1424--1433), ``Hahn supports
and accumulating action sets'' (1515--1524), ``Summation compatible with
refined inversion'' (1536--1543) and ``A staged Lean formalization target''
(1545--1562), which this article does not cite individually.
\end{ednote}

\section{Weights, excess degree, and the main statements}
\label{sec:statements}

\subsection{The exact type being measured}

Throughout, a weight means a sequence
\[
 N_0=1,\qquad N_n\in(0,\infty)\quad(n\ge1).
\]
It need not be increasing, log-convex, weakly log-convex, or stable under
a shift of its index. Let
\begin{equation}
 f(z)=z+\sum_{n\ge1}f_{n+1}z^{n+1}\in\C[[z]]
 \label{eq:tangent}
\end{equation}
be tangent to the identity. Its \emph{excess-degree type} is
\begin{equation}
 \Type(f)=\limsup_{n\to\infty}
 \left(\frac{|f_{n+1}|}{N_n}\right)^{1/n}\in[0,\infty].
 \label{eq:type}
\end{equation}
The excess degree of $z^{n+1}$ is $n$. This indexing is natural because
excess degrees add in the inverse coefficient formula. We write
$\inv f$ for the compositional inverse and reserve $1/f$ for a
multiplicative reciprocal.

To translate the usual coefficient weight $M_m$, put $N_n=M_{n+1}$ for
$n\ge1$. Replacing the exponent $1/n$ in \eqref{eq:type} by $1/(n+1)$
does not change the extended-real limsup. Thus the convention measures
exactly the same type as the repository's ordinary-degree convention,
with the harmless finite linear coefficient omitted.

For $A>0$, define the tangent dilation
\begin{equation}
 (\mathcal D_A f)(z)=A f(z/A).
 \label{eq:dilation}
\end{equation}
Then
\begin{equation}
 \Type(\mathcal D_A f)=A^{-1}\Type(f),\qquad
 \inv{(\mathcal D_A f)}=\mathcal D_A(\inv f).
 \label{eq:dilationproperties}
\end{equation}
If $\Type(f)<A$, there is a finite $C>0$ such that
\begin{equation}
 |f_{n+1}|\le C A^nN_n\qquad(n\ge1).
 \label{eq:coefficientbound}
\end{equation}
This follows directly from the definition of limsup, increasing $C$ to
cover the finitely many exceptional coefficients.

\begin{definition}[Envelope and distortion]
The \emph{supermultiplicative envelope} of $N$ is $S_0=1$ and
\begin{equation}
 S_n=\max_{\substack{k\ge1,\ j_1,\ldots,j_k\ge1\\
                         j_1+\cdots+j_k=n}}
             \prod_{i=1}^{k}N_{j_i}\qquad(n\ge1).
 \label{eq:envelope}
\end{equation}
Define
\begin{equation}
 \Delta(N)=\limsup_{n\to\infty}(S_n/N_n)^{1/n}\in[1,\infty].
 \label{eq:delta}
\end{equation}
We call $N$ \emph{admissible} if
\begin{equation}
 N_n^{1/n}\longrightarrow\infty,
 \qquad \log(S_n/N_n)=o(n).
 \label{eq:admissible}
\end{equation}
\end{definition}

There are only finitely many ordered compositions of each positive
integer, so \eqref{eq:envelope} is a finite maximum even for arbitrary
real weights. Positivity and finite support at each coefficient are all
that is required for the formal manipulations.

\subsection{The classification and its quantitative refinement}

\begin{theorem}[Complete zero-loss classification]
\label{thm:classificationintro}
For every positive weight $N$, the following are equivalent.
\begin{enumerate}[label=\textup{(\roman*)}]
\item Every tangent-to-identity scalar formal series satisfies
$\Type(\inv f)=\Type(f)$, including types zero and infinity.
\item $N$ is admissible in the sense of \eqref{eq:admissible}.
\item There exists a positive supermultiplicative sequence $V$, with
$V_0=1$, $V_n^{1/n}\to\infty$, and
$\log(V_n/N_n)=o(n)$.
\end{enumerate}
\end{theorem}

The equivalence is proved in Section~\ref{sec:classification}. The next
statement quantifies how it fails. Let $\mathcal I(N)$ be the infimum of
all finite $D\ge1$ such that
\begin{equation}
 \Type(\inv f)\le D\Type(f)
 \quad\hbox{for every tangent-to-identity } f.
 \label{eq:universalD}
\end{equation}
The infimum of the empty set is infinity; for finite $D$ we use
$D\cdot0=0$ and $D\cdot\infty=\infty$.

\begin{theorem}[Exact optimal distortion]
\label{thm:distortionintro}
For an arbitrary positive weight,
\begin{equation}
 \mathcal I(N)=
 \begin{cases}
 \Delta(N),&N_n^{1/n}\to\infty,\\
 \infty,&N_n^{1/n}\not\to\infty.
 \end{cases}
 \label{eq:optimaldistortion}
\end{equation}
When $\Delta(N)<\infty$ and the roots tend to infinity,
\begin{equation}
 \frac{\Type(f)}{\Delta(N)}\le\Type(\inv f)
             \le\Delta(N)\Type(f).
 \label{eq:twosideddistortion}
\end{equation}
The upper constant is attained by the explicit series
$f=z-\sum_{n\ge1}N_nz^{n+1}$, whose input type is one.
\end{theorem}

Thus there is a canonical obstruction, not just a list of sufficient
regularity conditions. For instance, a weight may have roots tending to
infinity yet have $\Delta(N)=2$, or $\Delta(N)=\infty$. Both possibilities
will be realized by elementary explicit formulas.

\subsection{The positive extremal series}

For a fixed $C>0$, put
\begin{equation}
 W_N(z)=\sum_{n\ge1}N_nz^{n+1},\qquad
 H_C(z)=\inv{(z-CW_N(z))}
       =z+\sum_{n\ge1}h_{n+1}(C)z^{n+1}.
 \label{eq:positiveinverse}
\end{equation}
It is equivalently the unique formal solution of
\begin{equation}
 H_C=z+C W_N(H_C).
 \label{eq:fixedpoint}
\end{equation}
All its coefficients are nonnegative. The following is the technical
core of the article.

\begin{theorem}[Envelope asymptotic]
\label{thm:envelopeintro}
If $N_n^{1/n}\to\infty$, then for every fixed $C>0$,
\begin{equation}
 \boxed{\displaystyle
 \lim_{n\to\infty}\left(\frac{h_{n+1}(C)}{S_n}\right)^{1/n}=1.}
 \label{eq:central}
\end{equation}
No bound on $S_n/N_n$ is assumed.
\end{theorem}

The assertion is about an exponential scale, not an asymptotic equivalent
$h_{n+1}(C)\sim S_n$. The residual ratio can grow or decay
subexponentially and may oscillate. Both the uniformity in the
coefficient index and the fact that $C$ may be smaller than one matter
in the proof.

\section{Envelope algebra and finite inverse coefficients}
\label{sec:envelope}

\subsection{A least supermultiplicative majorant}

\begin{lemma}[Envelope calculus]
\label{lem:envelopecalculus}
For every positive $N$, its envelope satisfies
\begin{equation}
 S_n\ge N_n,\qquad S_{a+b}\ge S_aS_b\quad(a,b\ge0),
 \qquad
 S_n=\max_{1\le j\le n}N_jS_{n-j}.
 \label{eq:envelopecalculus}
\end{equation}
It is the least pointwise supermultiplicative majorant of $N$ with zeroth
term one. It is unchanged if ordered compositions in
\eqref{eq:envelope} are replaced by unordered integer partitions.
\end{lemma}
\begin{proof}
The one-part composition proves $S_n\ge N_n$. Concatenating maximizing
compositions of $a$ and $b$ proves supermultiplicativity. Removing the
first part of a composition, and conversely adjoining a first part,
gives the recurrence. If $V\ge N$ and $V_{a+b}\ge V_aV_b$, repeated
multiplication gives
$V_n\ge\prod_iV_{j_i}\ge\prod_iN_{j_i}$ for every composition of $n$;
hence $V_n\ge S_n$. Finally, the product depends on the multiplicities of
parts, not their order.
\end{proof}

Writing $s_n=\log S_n$ and $a_n=\log N_n$ converts the recurrence to
\begin{equation}
 s_0=0,\qquad s_n=\max_{1\le j\le n}(a_j+s_{n-j}).
 \label{eq:maxplus}
\end{equation}
This is a finite max-plus closure. The obstruction to zero loss is the
excess $s_n-a_n$ on a linear scale. It is not the second difference of
$a_n$, nor a bound on a single adjacent quotient.

\subsection{Formal inversion and a residue calculation}

Every series of the form \eqref{eq:tangent} has a unique formal inverse:
the coefficient of degree $m$ in $f(g(z))=z$ determines $g_m$ using only
earlier coefficients. Associativity then gives a two-sided inverse.
All operations up to a prescribed degree use finitely many terms.

For completeness, the needed Lagrange identity can be obtained directly
from formal residues. If $g=\inv f$ and $m\ge1$, substitution $z=f(t)$
in a residue gives
\begin{align*}
 [z^m]g(z)
 &=\Res_z g(z)z^{-m-1}\,\dd z\\
 &=\Res_t t f(t)^{-m-1}f'(t)\,\dd t
  =\frac1m\Res_t f(t)^{-m}\,\dd t.
\end{align*}
The last equality follows because the residue of a formal derivative
vanishes, applied to $\dd(t f(t)^{-m})/\dd t$. Formal substitution is
legitimate since $f(t)=t$ times a unit. Therefore
\begin{equation}
 [z^m]g(z)=\frac1m[t^{m-1}]\left(\frac{t}{f(t)}\right)^m.
 \label{eq:lagrange}
\end{equation}
This is the classical formula used below \cite{gessel}.

\begin{lemma}[Positive coefficient formula and absolute majorant]
\label{lem:lagrange}
For \eqref{eq:positiveinverse} and $n\ge1$,
\begin{equation}
 h_{n+1}(C)=\frac1{n+1}\sum_{k=1}^{n}\binom{n+k}{k}C^k
 \sum_{\substack{j_1+\cdots+j_k=n\\j_i\ge1}}
                         \prod_{i=1}^{k}N_{j_i}.
 \label{eq:orderedformula}
\end{equation}
Equivalently, summing over finitely supported multiplicity vectors
$\boldsymbol m=(m_1,m_2,\ldots)$ with $\sum j m_j=n$ and
$k=\sum m_j$,
\begin{equation}
 h_{n+1}(C)=
 \sum_{\boldsymbol m}
 \frac{(n+k)!}{(n+1)!\prod_jm_j!}
 C^k\prod_j N_j^{m_j}.
 \label{eq:partitionformula}
\end{equation}
If $|f_{n+1}|\le C N_n$ for all $n$, then
\begin{equation}
 |[z^{n+1}]\inv f|\le h_{n+1}(C).
 \label{eq:absoluteinverse}
\end{equation}
\end{lemma}
\begin{proof}
In \eqref{eq:lagrange}, take $m=n+1$ and expand
\[
 \left(1-C\sum_{j\ge1}N_jt^j\right)^{-n-1}
 =\sum_{k\ge0}\binom{n+k}{k}C^k
                  \left(\sum_{j\ge1}N_jt^j\right)^k.
\]
Only $1\le k\le n$ contributes to degree $n$. Grouping the $k!/
\prod_jm_j!$ orderings of a fixed partition gives
\eqref{eq:partitionformula}. For a signed or complex input, the same
finite expansion and the triangle inequality give
\eqref{eq:absoluteinverse}.
\end{proof}

\begin{remark}[A coefficient lower bound]
\label{rem:lowerpartition}
For any ordered composition of $n$ with $k$ parts,
\begin{equation}
 h_{n+1}(C)\ge C^k\prod_iN_{j_i},
 \label{eq:lowerpartition}
\end{equation}
because $\binom{n+k}{k}/(n+1)\ge1$. In particular,
$h_{n+1}(1)\ge S_n$. A maximizing composition need not be unique; the
next lemma controls the lengths of all maximizers.
\end{remark}

\section{Sparse products, dense suppression, and the central limit}
\label{sec:centralproof}

\subsection{Maximizing partitions have sublinear length}

\begin{lemma}[Sparse maximizers]
\label{lem:sparsemax}
Assume $N_n^{1/n}\to\infty$. For every $\eta>0$, all compositions
attaining $S_n$ have fewer than $\eta n$ parts once $n$ is sufficiently
large. Equivalently, the maximum number of parts in a maximizing
composition is $o(n)$.
\end{lemma}
\begin{proof}
It suffices to consider $0<\eta\le1$. Set $L=\lceil2/\eta\rceil$ and
\[
 B=\max\left(1,\max_{1\le j\le L}N_j^{1/j}\right).
\]
Suppose a composition of $n$ has $k\ge\eta n$ parts. At most
$n/(L+1)<\eta n/2$ parts exceed $L$. Thus more than $\eta n/2$ parts
are at most $L$. Let $\ell$ be their total mass. Then
$\ell\ge\eta n/2$, and their product is at most $B^\ell$.

Since $N_\ell^{1/\ell}\to\infty$, for all sufficiently large $n$ we
have $N_\ell>B^\ell$ whenever $\ell\ge\eta n/2$. Replace those small
parts by the single part $\ell$. This preserves the sum $n$ and strictly
increases the product. Consequently the original composition was not
maximizing. The threshold is independent of the chosen composition.
\end{proof}

This argument does not assume monotone roots. The full root limit makes
the replacement by one part valid uniformly beyond a large index. A deep
dip could prevent that particular one-part replacement. A strengthened
version, using a whole maximizing subpartition instead, is given at the
end of this section.

\subsection{A finite upper certificate}

Define for $1\le k\le n$
\begin{equation}
 A_{n,k}=\frac1{n+1}\binom{n+k}{k}\binom{n-1}{k-1}.
 \label{eq:Ank}
\end{equation}
The second binomial counts ordered compositions of $n$ with $k$ positive
parts. For $0<\eta\le1/2$ and $C>0$, set
\begin{align}
 P_{n,\eta}(C)&=\sum_{1\le k\le\eta n}A_{n,k}C^k,
 &Q_{n,\eta}(C)&=\sum_{\eta n<k\le n}A_{n,k}C^k.
 \label{eq:PQ}
\end{align}
A sum with no indices is zero.

\begin{theorem}[Exact finite sparse--dense bound]
\label{thm:finitecertificate}
Fix any positive weight, $n\ge1$, $C>0$, and $0<\eta\le1/2$.
Let $L=\lceil2/\eta\rceil$ and choose $B\ge1$ satisfying
$N_j\le B^j$ for $1\le j\le\min(L,n)$. Put
\begin{equation}
 R_{n,\eta,B}=
 \max_{\lceil\eta n/2\rceil\le\ell\le n}\frac{B^\ell}{S_\ell}.
 \label{eq:Rtail}
\end{equation}
Then
\begin{equation}
 \boxed{\displaystyle
 h_{n+1}(C)\le S_n
 \bigl(P_{n,\eta}(C)+Q_{n,\eta}(C)R_{n,\eta,B}\bigr).}
 \label{eq:finitecertificate}
\end{equation}
No limiting hypothesis on $N$ is needed for this inequality.
\end{theorem}
\begin{proof}
For compositions with $k\le\eta n$, each product is at most $S_n$;
these give the first term. If $k>\eta n$, separate the parts at most
$L$. Their number exceeds $\eta n/2$, as in
Lemma~\ref{lem:sparsemax}, and their total mass $\ell$ satisfies
$\ell\ge\lceil\eta n/2\rceil$. Their product is at most $B^\ell$.
The product of the remaining parts is at most $S_{n-\ell}$, including
the empty product $S_0=1$. Since
$S_\ell S_{n-\ell}\le S_n$, the whole product divided by $S_n$ is at
most $B^\ell/S_\ell\le R_{n,\eta,B}$. Count compositions and apply
\eqref{eq:orderedformula}.
\end{proof}

All entries of \eqref{eq:finitecertificate} involve only the finite
prefix $N_1,\ldots,N_n$. For rational weights an integer $B$ can be
chosen by exact root comparisons, so this is an executable rational
inequality rather than a certificate requiring floating-point logarithms.

\subsection{Subexponential sparse entropy}

Let $\ent(t)=-t\log t-(1-t)\log(1-t)$ for $0<t<1$, extended by zero
at the endpoints. The elementary binomial entropy inequality is
\begin{equation}
 \binom{r}{s}\le\exp\bigl(r\ent(s/r)\bigr).
 \label{eq:binomialentropy}
\end{equation}
For $0<s<r$, it follows by taking the $s$th term in
$1=(s/r+(1-s/r))^r$ and rearranging; the endpoint cases are immediate.

\begin{lemma}[Counting bounds]
\label{lem:counting}
For $0<\eta\le1/2$,
\begin{equation}
 P_{n,\eta}(C)\le n\exp\bigl(n E_C(\eta)\bigr),
 \quad
 E_C(\eta)=2\ent(\eta/2)+\ent(\eta)+\eta\lp C,
 \label{eq:sparseentropy}
\end{equation}
where $E_C(\eta)\to0$ as $\eta\downarrow0$. Also
\begin{equation}
 Q_{n,\eta}(C)\le
 n\bigl(8\max(1,C)\bigr)^n.
 \label{eq:densecount}
\end{equation}
\end{lemma}
\begin{proof}
For $k\le\eta n$, use
$\binom{n+k}{k}\le\binom{2n}{k}$ and
$\binom{n-1}{k-1}=(k/n)\binom nk\le\binom nk$.
The function $\ent$ is increasing on $[0,1/2]$, so
\eqref{eq:binomialentropy} bounds the binomials by
$\exp(2n\ent(\eta/2))$ and $\exp(n\ent(\eta))$, respectively.
Also $C^k\le\exp(n\eta\lp C)$, and there are at most $n$ summands.
For the second assertion use the cruder bounds
$\binom{n+k}{k}\le2^{2n}$,
$\binom{n-1}{k-1}\le2^n$, and
$C^k\le\max(1,C)^n$.
\end{proof}

The sparse count has arbitrarily small exponential rate. The dense count
can have a fixed exponential rate, but its products incur a stronger
penalty from the weight. This separation is the reason no smoothness of
the weight profile is required.

\begin{lemma}[Dense suppression]
\label{lem:densesuppression}
If $N_n^{1/n}\to\infty$, then for fixed $\eta>0$ and fixed $B\ge1$,
\begin{equation}
 R_{n,\eta,B}^{1/n}\longrightarrow0.
 \label{eq:densesuppression}
\end{equation}
\end{lemma}
\begin{proof}
We have $S_\ell\ge N_\ell$. Given $K>B$, eventually
$S_\ell^{1/\ell}\ge K$ for every
$\ell\ge\lceil\eta n/2\rceil$. Therefore
\[
 R_{n,\eta,B}\le(B/K)^{\eta n/2},\qquad
 \limsup_nR_{n,\eta,B}^{1/n}\le(B/K)^{\eta/2}.
\]
Let $K\to\infty$.
\end{proof}

\subsection{Proof of the envelope asymptotic}

\begin{proof}[Proof of Theorem~\ref{thm:envelopeintro}]
For each $n$, select a maximizing composition with $k_n$ parts. By
Lemma~\ref{lem:sparsemax}, $k_n/n\to0$. The lower bound
\eqref{eq:lowerpartition} gives
\[
 \left(h_{n+1}(C)/S_n\right)^{1/n}\ge C^{k_n/n}\longrightarrow1.
\]
This works for every fixed $C>0$, including $C<1$.

For the upper bound, fix $0<\eta\le1/2$, let
$L=\lceil2/\eta\rceil$, and choose once and for all
$B\ge\max(1,\max_{1\le j\le L}N_j^{1/j})$. Combining
Theorem~\ref{thm:finitecertificate} and Lemmas~\ref{lem:counting} and
\ref{lem:densesuppression} gives
\begin{align*}
 \limsup_n\left(h_{n+1}(C)/S_n\right)^{1/n}
 &\le\max\left(\ee^{E_C(\eta)},0\right)
 =\ee^{E_C(\eta)}.
\end{align*}
Here the $n$th root of the dense term tends to zero despite its fixed
exponential counting factor. Letting $\eta\downarrow0$ gives an upper
bound of one, and proves \eqref{eq:central}.
\end{proof}

\begin{corollary}[Type of the extremal inverse]
\label{cor:extremaltype}
If $N_n^{1/n}\to\infty$, then for every fixed $C>0$,
\begin{equation}
 \Type(H_C)=\Delta(N),\qquad \Type(z-CW_N)=1.
 \label{eq:extremaltype}
\end{equation}
\end{corollary}
\begin{proof}
Multiply $(h_{n+1}(C)/S_n)^{1/n}\to1$ by
$(S_n/N_n)^{1/n}$ and take limsups, allowing infinity. For the input,
$C^{1/n}\to1$.
\end{proof}

\subsection{A strengthening for unbounded root profiles}

\begin{proposition}[The envelope needs only unbounded primitive roots]
\label{prop:unboundedroots}
For a positive weight $N$, the following are equivalent:
\[
 \sup_{n\ge1}N_n^{1/n}=\infty,
 \qquad S_n^{1/n}\longrightarrow\infty.
\]
Under either condition the conclusion \eqref{eq:central} holds for every
fixed $C>0$, even if $N_n^{1/n}$ itself does not tend to infinity.
\end{proposition}
\begin{proof}
If the primitive roots are bounded by $K$, every composition product is
at most $K^n$, and $S_n\le K^n$. Conversely, fix $m\ge1$ and write
$n=qm+r$ with $0\le r<m$. Supermultiplicativity gives
$S_n\ge N_m^qS_r$, so
$\liminf_n S_n^{1/n}\ge N_m^{1/m}$. Taking arbitrarily large primitive
roots proves the equivalence.

In the proof of Lemma~\ref{lem:sparsemax}, the small parts have mass
$\ell\ge\eta n/2$ and product at most $B^\ell$. Now
$S_\ell>B^\ell$ for all such $\ell$ once $n$ is large. Replace those
small parts by any composition attaining $S_\ell$, rather than by the
single part $\ell$. This again strictly increases the product, so all
maximizing compositions still have sublinear length. The lower bound
for $h_{n+1}(C)/S_n$ follows unchanged. The finite upper certificate
already holds for arbitrary weights, and the dense-suppression proof
only needs $S_\ell^{1/\ell}\to\infty$. The entropy argument therefore
proves \eqref{eq:central} under the weaker hypothesis.
\end{proof}

The full limit $N_n^{1/n}\to\infty$ remains necessary for universal
zero loss, by the polynomial obstruction below. The strengthening
concerns the extremal envelope asymptotic, not an exception to that
classification. If the primitive roots are unbounded but have a bounded
subsequence, $S_n^{1/n}\to\infty$ forces $\Delta(N)=\infty$.

\section{Complete classification, optimal loss, and canonical repair}
\label{sec:classification}

\subsection{Why superexponential roots are necessary}

\begin{lemma}[The polynomial obstruction]
\label{lem:polynomial}
If a finite $D$ satisfies \eqref{eq:universalD}, then
$N_n^{1/n}\to\infty$.
\end{lemma}
\begin{proof}
The polynomial $p(z)=z-z^2$ has type zero for every weight. Its inverse is
\[
 \inv p(z)=\frac{1-\sqrt{1-4z}}2
 =\sum_{n\ge0}\frac1{n+1}\binom{2n}{n}z^{n+1}.
\]
The central binomial coefficient satisfies
$4^n/(2n+1)\le\binom{2n}{n}\le4^n$: it is the largest of the
$2n+1$ binomial coefficients whose sum is $4^n$. Consequently the $n$th
root of the coefficient of $z^{n+1}$ tends to four. If
$N_n^{1/n}$ does not tend to infinity, some subsequence has
$N_n^{1/n}\le B<\infty$, and
$\Type(\inv p)\ge4/B>0$. This contradicts
$\Type(\inv p)\le D\Type(p)=0$.
\end{proof}

The obstruction concerns type zero, not just the growth of a generic
divergent series. It explains why a criterion based solely on the envelope
defect would be incomplete. For example, $N_n=1$ is already
supermultiplicative and has $S_n=N_n$, but it cannot preserve the type of
an inverse polynomial.

\subsection{The exact distortion constant}

\begin{proof}[Proof of Theorem~\ref{thm:distortionintro}]
If the roots do not tend to infinity,
Lemma~\ref{lem:polynomial} rules out every finite $D$, as claimed.
Assume henceforth that $N_n^{1/n}\to\infty$.

First suppose $\Delta=\Delta(N)<\infty$ and $\Type(f)<\infty$.
For each $A>\Type(f)$, choose $C$ in
\eqref{eq:coefficientbound}. The coefficients of $\mathcal D_A f$
are bounded in absolute value by $CN_n$. Thus
Lemma~\ref{lem:lagrange} and Corollary~\ref{cor:extremaltype} give
\[
 A^{-1}\Type(\inv f)
 =\Type\bigl(\inv{(\mathcal D_A f)}\bigr)
 \le\Type(H_C)=\Delta.
\]
Let $A\downarrow\Type(f)$. In particular a type-zero series has a
type-zero inverse. The upper bound is automatic when $\Type(f)=\infty$.
It therefore holds for every $f$. Applying it to $\inv f$ proves the
lower bound in \eqref{eq:twosideddistortion}; when $\Type(f)=\infty$,
a finite inverse type would contradict that same inequality.

The input $z-W_N$ has type one and its inverse has type $\Delta$ by
Corollary~\ref{cor:extremaltype}. Hence no smaller constant works.
If $\Delta=\infty$, this same input rules out every finite constant.
This proves the formula in all cases.
\end{proof}

The result controls a limsup type, not the ratio of individual signed
inverse coefficients to their inputs. Individual coefficients may cancel
completely. The positive extremal series removes cancellation and
identifies the optimal universal obstruction.

\subsection{Zero loss and equivalence of weights}

Two positive weights will be called \emph{subexponentially equivalent}
when
\begin{equation}
 \log(V_n/N_n)=o(n).
 \label{eq:weightequivalence}
\end{equation}
They define the same exact type for every coefficient sequence: multiplying
each normalized $n$th root by $(N_n/V_n)^{1/n}\to1$ preserves every
extended-real limsup.

\begin{proof}[Proof of Theorem~\ref{thm:classificationintro}]
Assume (i). The polynomial obstruction gives superexponential roots.
The input $z-W_N$ and Corollary~\ref{cor:extremaltype} then give
$\Delta(N)=1$. Since $S_n/N_n\ge1$, this is equivalent to
$\log(S_n/N_n)=o(n)$, proving (ii).

If (ii) holds, Theorem~\ref{thm:distortionintro} with $\Delta=1$
proves (i), and $V=S$ proves (iii).

Finally suppose (iii). Since $V$ is supermultiplicative, its own envelope
is exactly $V$. Its roots tend to infinity, so it is admissible and the
already proved implication (ii)$\Rightarrow$(i) applies to $V$.
Subexponential equivalence identifies the $V$-type and $N$-type of every
series and its inverse. Hence (i) also holds for $N$.
\end{proof}

\begin{corollary}[Stability under subexponential perturbations]
\label{cor:perturbations}
If $N$ is admissible and $\log(V_n/N_n)=o(n)$, then $V$ is admissible.
In particular, altering finitely many weights by arbitrary positive
factors does not affect admissibility or any exact type.
\end{corollary}
\begin{proof}
All exact types coincide for the two weights. Transfer the universal
inversion identity and apply Theorem~\ref{thm:classificationintro}.
\end{proof}

This proof deliberately uses the characterization rather than multiplying
pointwise $\exp(o(n))$ errors over an arbitrary partition. A partition
could contain linearly many small parts; such an unqualified multiplication
would lose the very precision at issue.

\subsection{A minimal admissible repair}

\begin{theorem}[Canonical envelope repair]
\label{thm:repair}
Suppose $N_n^{1/n}\to\infty$. Its envelope $S$ is admissible.
For any admissible pointwise majorant $V\ge N$,
\begin{equation}
 \limsup_{n\to\infty}(S_n/V_n)^{1/n}\le1.
 \label{eq:minimalrepair}
\end{equation}
Thus $S$ is the least supermultiplicative majorant pointwise, and is a
minimal admissible repair up to subexponential domination.
\end{theorem}
\begin{proof}
The sequence $S$ is supermultiplicative and has roots tending to infinity,
so its own envelope is $S$ and it is admissible. Let $S^{V}$ denote the
envelope of $V$. Monotonicity of the maximum-product construction gives
$S_n\le S^{V}_n$. Since $V$ is admissible,
$(S^{V}_n/V_n)^{1/n}\to1$, proving \eqref{eq:minimalrepair}.
\end{proof}

\begin{remark}[The precise minimality claim]
We do not claim that $S$ is the pointwise smallest among all admissible
majorants. An admissible weight may have a finite or subexponential defect
from exact supermultiplicativity. The pointwise minimality applies to
supermultiplicative majorants; \eqref{eq:minimalrepair} is the separate
asymptotic statement for admissible ones.
\end{remark}

\subsection{Recovery of the log-convex theorem}

\begin{proposition}[Log-convex weights need no eventual repair]
\label{prop:logconvex}
Let $M_0=1$ and let $M$ be positive and log-convex, with
$M_m^{1/m}\to\infty$. Set $N_0=1$ and $N_n=M_{n+1}$ for $n\ge1$.
Then its envelope satisfies $S_n=N_n$ for every sufficiently large $n$.
In particular $N$ is admissible.
\end{proposition}
\begin{proof}
Log-convexity says the quotients $M_j/M_{j-1}$ are nondecreasing.
For integers $a\ge b\ge2$, transferring one unit from $b$ to $a$
therefore gives
$M_aM_b\le M_{a+1}M_{b-1}$ whenever $b\ge3$.
Repeatedly transfer until all but one indices equal two. For
$j_1+\cdots+j_k=n$, with $j_i\ge1$, this gives
\begin{equation}
 \prod_{i=1}^kM_{j_i+1}\le M_{n-k+2}M_2^{k-1}.
 \label{eq:logconvexconcentration}
\end{equation}
Since $M_0=1$, convexity of $m\mapsto\log M_m$ implies
$M_r\le M_{n+1}^{r/(n+1)}$ for $0\le r\le n+1$. Thus the
right side of \eqref{eq:logconvexconcentration} is at most
\[
 M_{n+1}
 \left(\frac{M_2}{M_{n+1}^{1/(n+1)}}\right)^{k-1}.
\]
For all sufficiently large $n$ the parenthesis is at most one,
uniformly in $k$. Every partition product is then at most $M_{n+1}$,
and the one-part partition attains that value. Also
$N_n^{1/n}=\bigl(M_{n+1}^{1/(n+1)}\bigr)^{(n+1)/n}\to\infty$.
\end{proof}

This is the log-convex special case treated in the predecessor
\cite{repo-weight}, recovered here rather than imported as a new result.
It also explains why the one-step shift in the definition of $N$ is
important: the products in inversion involve $M_{j_i+1}$, not $M_{j_i}$.
The shifted weight itself need not satisfy a normalization-dependent
log-convex condition at its first few indices.

\section{Two divergent composition arguments and exact cancellation}
\label{sec:composition}

In this section $N$ is admissible. No convergence of a formal series in
a nonzero disk is assumed.

\subsection{The maximum bound}

For nonnegative-coefficient series, write $F\preceq G$ when every
coefficient of $F$ is at most the corresponding coefficient of $G$.
Substitution by series with zero constant coefficient preserves this
order. Put $B_C=z+CW_N$.

\begin{lemma}[A single inverse majorizes a double composition]
\label{lem:doublemajorant}
For every $C>0$,
\begin{equation}
 B_C\circ B_C\preceq H_{2C}.
 \label{eq:doublemajorant}
\end{equation}
\end{lemma}
\begin{proof}
The identity $H_{2C}=z+2C W_N(H_{2C})$ gives
$z\preceq B_C\preceq H_{2C}$. Consequently
\[
 B_C(B_C)=z+CW_N(z)+CW_N(B_C)
 \preceq z+2CW_N(H_{2C})=H_{2C}.
\]
\end{proof}

\begin{theorem}[Sharp composition law]
\label{thm:composition}
For tangent-to-identity formal series $f,g$,
\begin{equation}
 \boxed{\displaystyle
 \Type(f\circ g)\le\max\{\Type(f),\Type(g)\}.}
 \label{eq:compositionmax}
\end{equation}
If $\Type(f)\ne\Type(g)$, equality holds, including when one type is
infinite.
\end{theorem}
\begin{proof}
If both types are finite, take
$A>\max\{\Type(f),\Type(g)\}$. After the common dilation
$\mathcal D_A$, choose one $C$ bounding the absolute nonlinear
coefficients of both inputs by $CN_n$. Their composite is bounded
coefficientwise in absolute value by $B_C\circ B_C$, hence by $H_{2C}$.
Admissibility and Corollary~\ref{cor:extremaltype} give type at most
one after dilation. Since dilation commutes with composition, the
undilated type is at most $A$. Let $A$ decrease to the maximum. If an
input type is infinite, the upper bound is automatic.

Suppose $\Type(f)>\Type(g)$. From
$f=(f\circ g)\circ\inv g$, the already proved upper bound and exact
inverse invariance give
\[
 \Type(f)\le\max\{\Type(f\circ g),\Type(g)\}.
\]
This forces $\Type(f\circ g)\ge\Type(f)$, also when $\Type(f)=\infty$.
The other strict inequality is handled by
$g=\inv f\circ(f\circ g)$.
\end{proof}

The estimate is not a product bound. Tangent dilations normalize both
inputs at the same exponential scale, while every finite composition
has only a subexponential overhead relative to the admissible envelope.
The inverse identity then turns an upper bound into an exact law whenever
the two scales differ.

\subsection{Every allowed cancellation actually occurs}

\begin{theorem}[Complete type-only cancellation spectrum]
\label{thm:cancellation}
Fix two types $s,t\in[0,\infty]$.
If $s\ne t$, every composite of tangent series of these types has type
$\max(s,t)$. If $s=t=T$, the set of possible composite types is exactly
$[0,T]$, with the natural extended-real interpretation for $T=\infty$.
\end{theorem}
\begin{proof}
The unequal case and the upper restriction in the equal case follow from
Theorem~\ref{thm:composition}. Every finite positive type $t$ is realized
by the positive series
\[
 f_t(z)=z+\sum_{n\ge1}t^nN_nz^{n+1}.
\]
The identity realizes zero, and
$z+\sum_{n\ge1}n^nN_nz^{n+1}$ realizes infinity.

Let $0\le S<T\le\infty$. Choose tangent $f$ of type $T$ and $h$ of
type $S$, and set $g=\inv f\circ h$. The unequal-type law and inverse
invariance imply $\Type(g)=T$. Yet $f\circ g=h$ has type $S$.
For $S=T>0$, choose $f=g$ with nonnegative nonlinear coefficients and
type $T$. Since $f(z)\succeq z$, we have $f\circ f\succeq f$, so
the composite has type at least $T$, and hence exactly $T$. This also
works at infinity. At $T=0$, the upper bound already forces zero.
\end{proof}

Thus there is no missing universal scalar formula in the equal-type case:
all values allowed by the upper bound occur. A more precise prediction
must use information beyond the two types, such as signs, coefficient
profiles, or a specified relation between the series.

\subsection{Type-zero coordinates need not be analytic}

\begin{corollary}[Type-zero coordinate invariance]
\label{cor:zerocoordinates}
If $h$ is tangent to the identity and $\Type(h)=0$, then for every
tangent $f$,
\begin{equation}
 \Type(h\circ f)=\Type(f\circ h)=\Type(f),
 \qquad
 \Type(\inv h\circ f\circ h)=\Type(f).
 \label{eq:zerocoordinates}
\end{equation}
\end{corollary}
\begin{proof}
Use the unequal-type law when $\Type(f)>0$, and the maximum bound when
it is zero. The inverse of $h$ also has type zero, so the conjugation
statement follows by applying the two identities successively.
\end{proof}

Every analytic tangent germ has type zero because its coefficients have
at most exponential growth while $N_n^{1/n}\to\infty$. But the
corollary is strictly more general. For example,
\begin{equation}
 h(z)=z+\sum_{n\ge1}\sqrt{N_n}\,z^{n+1}
 \label{eq:divergentzero}
\end{equation}
has radius of convergence zero, since
$(\sqrt{N_n})^{1/(n+1)}\to\infty$, whereas
\[
 \Type(h)=\limsup_n N_n^{-1/(2n)}=0.
\]
It is therefore a genuinely divergent type-zero coordinate change under
which every exact $N$-type is preserved.

\subsection{Nonunit linear coefficients}

For a formal series with zero constant coefficient and nonzero linear
coefficient, use the same definition \eqref{eq:type} on its nonlinear
part. Multiplying the whole series by a fixed nonzero scalar does not
change the type. Precomposing by $z\mapsto bz$ multiplies type by $|b|$.

\begin{corollary}[Scalar coordinate scaling]
\label{cor:scalarscaling}
Let $F'(0)=a\ne0$ and $G'(0)=b\ne0$. Then
\begin{align}
 \Type(\inv F)&=|a|^{-1}\Type(F),\label{eq:scalarinverse}\\
 \Type(F\circ G)&\le
       \max\{|b|\Type(F),\Type(G)\}.
       \label{eq:scalarcomposition}
\end{align}
Equality holds in \eqref{eq:scalarcomposition} whenever the two entries
of the maximum differ. If $h,k$ are invertible type-zero series with
$h'(0)=b$, then
\begin{equation}
 \Type(k\circ F\circ h)=|b|\Type(F).
 \label{eq:scalarcoordinates}
\end{equation}
\end{corollary}
\begin{proof}
Write $F=af$ and $G=bg$ with $f,g$ tangent. Then
$\inv F=(\inv f)(z/a)$, proving \eqref{eq:scalarinverse}. Also
\[
 F\circ G=ab\bigl(f_b\circ g\bigr),\qquad
 f_b(z)=b^{-1}f(bz),\qquad \Type(f_b)=|b|\Type(F).
\]
Apply Theorem~\ref{thm:composition}. Applying the result twice with a
type-zero outer or inner series proves \eqref{eq:scalarcoordinates};
when the surviving type is zero, the upper bound itself gives equality.
\end{proof}

\subsection{A filtered formal group}

Let $\mathcal G_{\le t}$ be the tangent formal series of type at most
$t$. Exact inversion and the maximum bound show that each
$\mathcal G_{\le t}$ is a group under composition. The function
\begin{equation}
 d(f,g)=\Type(f\circ\inv g)
 \label{eq:pseudodistance}
\end{equation}
is a right-invariant extended ultrapseudometric:
\[
 d(f,g)=d(g,f),\qquad
 d(f,h)\le\max\{d(f,g),d(g,h)\},\qquad
 d(f\circ k,g\circ k)=d(f,g).
\]
Indeed the first identity follows from inversion invariance, the second
from composition, and the third from cancellation of $k\circ\inv k$.
It is a pseudometric rather than a metric because nontrivial type-zero
series have zero distance from the identity. No left-invariance or
normality of the type-zero subgroup is asserted.

\section{Formal maps in several variables and Banach spaces}
\label{sec:banach}

\subsection{The coefficient norm must be specified}

Let $E$ be a real or complex Banach space. A formal map tangent to the
identity is a sequence of bounded symmetric multilinear maps
$A_m:E^m\to E$, written
\begin{equation}
 F(x)=x+\sum_{m\ge2}A_m(x,\ldots,x).
 \label{eq:banachmap}
\end{equation}
No convergence of the displayed series is required. The coefficient norm
is the multilinear operator norm
\begin{equation}
 \|A_m\|=\sup_{\|x_1\|,\ldots,\|x_m\|\le1}
                   \|A_m(x_1,\ldots,x_m)\|.
 \label{eq:multilinearnorm}
\end{equation}
Define
\begin{equation}
 \Type(F)=\limsup_{n\to\infty}
              \left(\frac{\|A_{n+1}\|}{N_n}\right)^{1/n}.
 \label{eq:banachtype}
\end{equation}
The norm in \eqref{eq:multilinearnorm} is not silently replaced by the
supremum of the diagonal polynomial on the unit ball. Polarization
comparisons can have degree-dependent constants; exact exponential types
must therefore keep the norm convention explicit.

For an $m$-linear map $L$, write
\[
 \Sym L(x_1,\ldots,x_m)=\frac1{m!}
 \sum_{\sigma\in\mathfrak S_m}L(x_{\sigma(1)},\ldots,x_{\sigma(m)}).
\]
Then $\|\Sym L\|\le\|L\|$. Normalized symmetrization, rather than an
unnormalized sum, is essential here.

\subsection{Scalar majorants survive multilinear contraction}

Suppose $G(x)=\sum_{j\ge1}G_j(x,\ldots,x)$ has zero constant term.
The degree-$m$ coefficient in $F\circ G$ is
\begin{equation}
 (F\circ G)_m=
 \sum_{k=1}^m\ \sum_{\substack{j_1+\cdots+j_k=m\\j_i\ge1}}
 \Sym\bigl(A_k(G_{j_1},\ldots,G_{j_k})\bigr),
 \label{eq:banachcomposition}
\end{equation}
where $A_1=\id$ for a tangent map. In each summand the $m$ arguments
are first divided into consecutive blocks of sizes $j_i$, then fed into
the corresponding $G_{j_i}$, then into $A_k$; the result is symmetrized.
On the diagonal, expansion of the polynomial composition gives exactly
this formula.

The norm of a summand is at most
\begin{equation}
 \|A_k\|\prod_{i=1}^k\|G_{j_i}\|.
 \label{eq:contractionnorm}
\end{equation}
There is no factor depending on $\dim E$ and no extra factor $m!$.
Thus scalar series with nonnegative coefficient bounds majorize all
composition coefficients in this convention.

\begin{lemma}[Formal inverse and its scalar majorant]
\label{lem:banachinverse}
A formal map \eqref{eq:banachmap} has a unique two-sided formal inverse
$G$ tangent to the identity. If $\|A_{n+1}\|\le CN_n$, then
\begin{equation}
 \|G_{n+1}\|\le h_{n+1}(C)\qquad(n\ge1).
 \label{eq:banachmajorant}
\end{equation}
\end{lemma}
\begin{proof}
Set $G_1=\id$. At degree $m\ge2$, the equation $F\circ G=\id$
requires
\begin{equation}
 G_m=-\sum_{k=2}^m\ \sum_{j_1+\cdots+j_k=m}
          \Sym\bigl(A_k(G_{j_1},\ldots,G_{j_k})\bigr).
 \label{eq:banachinverserecursion}
\end{equation}
Every $j_i$ in the sum is at most $m-1$, so the formula is a finite
recursion defining a bounded symmetric coefficient. Formal composition
is associative because every coefficient identity is an identity of
finite polynomial substitutions. The same recursion gives a right
inverse for $G$; associativity then makes it equal to $F$, proving that
the inverse is two-sided.

Taking norms in \eqref{eq:banachinverserecursion}, applying
\eqref{eq:contractionnorm}, and inducting on $m$ yields the scalar
coefficient recursion for $H_C=z+CW_N(H_C)$. This proves
\eqref{eq:banachmajorant}.
\end{proof}

\subsection{Exact tangent calculus without a dimension loss}

\begin{theorem}[Banach-space zero-loss calculus]
\label{thm:banach}
Let $N$ be admissible. For formal maps tangent to the identity on any
real or complex Banach space, with norm \eqref{eq:multilinearnorm},
\begin{align}
 \Type(\inv F)&=\Type(F),\label{eq:banachinverseexact}\\
 \Type(F\circ G)&\le\max\{\Type(F),\Type(G)\}.
 \label{eq:banachcompositionexact}
\end{align}
The composition inequality is an equality when the two types differ.
Type-zero tangent precomposition, postcomposition, and conjugation
preserve type.
\end{theorem}
\begin{proof}
Dilation $\mathcal D_A F(x)=A F(x/A)$ divides a homogeneous coefficient
of degree $n+1$ by $A^n$, for positive real $A$. If the type is finite,
normalize it with any larger $A$ and use
Lemma~\ref{lem:banachinverse} and $\Type(H_C)=1$ to obtain the inverse
upper bound. Let $A$ decrease to the type, and apply the bound to the
inverse to get equality, exactly as in the scalar proof. Infinite types
are handled by the same involution argument.

For composition, \eqref{eq:banachcomposition} transfers the scalar
majorant $B_C\circ B_C\preceq H_{2C}$ to the multilinear norms.
Dilation gives the maximum bound. Associativity, the inverse equality,
and the identities used in Theorem~\ref{thm:composition} give equality
for unequal types. The type-zero assertions follow in the same way.
\end{proof}

The theorem concerns formal coefficient norms. It is not an analytic
inverse-function theorem on an open ball, nor does it assert that a
formal map is the Taylor expansion of a smooth map in a prescribed class.
The scalar majorant argument removes a combinatorial exponential loss;
it does not supply convergence where none was assumed.

\subsection{A coefficient-sum convention in finite dimensions}

There is a second useful norm convention for direct computation. For a
homogeneous polynomial map $P:\C^d\to\C^d$ of degree $m$, put
\begin{equation}
 \|P\|_{\mathrm{coeff}}=
 \max_{1\le r\le d}\sum_{|\alpha|=m}|[x^\alpha]P_r(x)|.
 \label{eq:coefficientnorm}
\end{equation}
Multiplying scalar polynomials multiplies coefficient-sum majorants.
For every monomial of the outer component, replacing all input variables
by one common scalar majorant bounds its entire substitution. Summing the
outer coefficients and taking the maximum over components therefore
proves the same scalar composition bound as
\eqref{eq:contractionnorm}. The inverse recursion and
Theorem~\ref{thm:banach} consequently hold for the type defined using
\eqref{eq:coefficientnorm} as well.

This is the convention used by the finite bivariate checks. We do not
identify it with the multilinear operator norm or transfer exact types
between norm conventions by an unstated degree-dependent equivalence.
The two theorems follow from their respective scalar majorization
properties.

\subsection{General derivatives: sharp bounds, not a single scalar law}

\begin{theorem}[Sharp Jacobian bounds]
\label{thm:jacobian}
Let $N$ be admissible and let a formal map $F$ on a nonzero Banach space
have bounded invertible linear coefficient $A$. Then
\begin{equation}
 \boxed{\displaystyle
 \frac{\Type(F)}{\|A\|}\le\Type(\inv F)
                         \le\|A^{-1}\|\Type(F).}
 \label{eq:jacobianbounds}
\end{equation}
Both endpoints are sharp even in dimension two. In general the matrix
$A$ and the scalar value $\Type(F)$ do not determine
$\Type(\inv F)$.
\end{theorem}
\begin{proof}
Multiplication of all output coefficients by a fixed bounded invertible
operator does not change type, since its norm and inverse norm are
constant factors whose $n$th roots tend to one. Precomposition by a
nonzero bounded linear operator $B$ gives
\[
 \|P_{n+1}(B,\ldots,B)\|\le\|P_{n+1}\|\|B\|^{n+1},
\]
and hence multiplies type by at most $\|B\|$.

Write $F=A\circ f$ with $f$ tangent. Then
$\inv F=\inv f\circ A^{-1}$, so
\[
 \Type(\inv F)\le\|A^{-1}\|\Type(\inv f)
                   =\|A^{-1}\|\Type(f)
                   =\|A^{-1}\|\Type(F).
\]
Apply this inequality to $\inv F$, whose linear coefficient is
$A^{-1}$, to obtain the lower bound.

For sharpness, take $E=\C^2$ with the maximum norm and
$A=\operatorname{diag}(a,b)$, where $0<a<b$. Let $\varphi$ be a
scalar tangent series of finite positive type $T$, and define
\[
 F_1(x,y)=(a\varphi(x),by),\qquad
 F_2(x,y)=(ax,b\varphi(y)).
\]
Both have derivative $A$ and type $T$. Their inverses are
\[
 \inv F_1(x,y)=(\inv\varphi(x/a),y/b),\qquad
 \inv F_2(x,y)=(x/a,\inv\varphi(y/b)),
\]
and have types $T/a$ and $T/b$, respectively. Since
$\|A^{-1}\|=1/a$ and $\|A\|=b$, the endpoints are attained with the
same derivative and the same scalar input type.
\end{proof}

The obstruction is directional: the growing nonlinear coefficients can
be placed in different eigendirections. An exact general-Jacobian theory
must therefore retain directional or multihomogeneous growth data, not
merely improve a combinatorial estimate in a scalar norm.

\section{Explicit weights separating the possible behaviors}
\label{sec:examples}

The examples in this section require no undecided asymptotic input. Their
envelopes are computed exactly. Together they show that the classification
is neither a disguised log-convexity condition nor a consequence of root
growth alone.

\subsection{A staircase weight beyond every log-convex approximation}

\begin{theorem}[Non-convex admissibility]
\label{thm:staircase}
The weight
\begin{equation}
 N_0=1,\qquad N_n=2^{\,n\lfloor\log_2n\rfloor}\quad(n\ge1)
 \label{eq:staircase}
\end{equation}
is supermultiplicative and admissible. There is no positive eventually
log-convex sequence $V$ satisfying $\log(V_n/N_n)=o(n)$.
\end{theorem}
\begin{proof}
The roots $r_n=N_n^{1/n}=2^{\lfloor\log_2n\rfloor}$ are
nondecreasing and tend to infinity. For any composition of $n$,
\[
 \prod_iN_{j_i}=\prod_i r_{j_i}^{j_i}
                  \le r_n^{\sum_i j_i}=N_n.
\]
Thus $S_n=N_n$, proving supermultiplicativity and admissibility.

Write $a_n=\log_2N_n=n\lfloor\log_2n\rfloor$. At $n=2^m$ with
$m\ge1$, direct calculation gives
\begin{equation}
 2a_n-a_{n-1}-a_{n+1}=n-1.
 \label{eq:staircasedefect}
\end{equation}
Indeed $a_{n-1}=(n-1)(m-1)$, $a_n=nm$, and
$a_{n+1}=(n+1)m$. If $\log_2V_n=a_n+e_n$ with $e_n=o(n)$,
then $2e_n-e_{n-1}-e_{n+1}=o(n)$ along these indices. The convexity
defect for $\log_2V$ is therefore $n-1+o(n)>0$, contradicting eventual
log-convexity, which requires that defect to be at most zero.
\end{proof}

This disproves the necessity of the log-convex-envelope route suggested
as one possible approach in the repository's arbitrary-weight question
\cite{repo-weight}. Subexponential equivalence to a supermultiplicative
sequence is necessary and sufficient; subexponential equivalence to a
log-convex sequence is strictly more restrictive.

\subsection{Admissibility without asymptotically nondecreasing roots}

\begin{proposition}[Persistent downward root jumps]
\label{prop:alternating}
The weight
\begin{equation}
 N_0=1,\qquad N_n=2^{n^2+(-1)^n n}\quad(n\ge1)
 \label{eq:alternating}
\end{equation}
is supermultiplicative and admissible. At every even $n$,
\begin{equation}
 \frac{N_{n+1}^{1/(n+1)}}{N_n^{1/n}}=\frac12.
 \label{eq:rootdrop}
\end{equation}
It is not subexponentially equivalent to a sequence whose roots are
eventually nondecreasing.
\end{proposition}
\begin{proof}
Let $a_n=n^2+(-1)^n n$. For positive $u,v$, the difference
$a_{u+v}-a_u-a_v$ equals $2uv$ when both are even, equals
$2u(v-1)$ when $u$ is even and $v$ odd, and equals $2uv+2(u+v)$
when both are odd. The remaining mixed case is symmetric. Each difference
is nonnegative, so $N$ is supermultiplicative. Its roots are
$2^{n+(-1)^n}\to\infty$, proving admissibility and
\eqref{eq:rootdrop}.

If $V$ were subexponentially equivalent to $N$, then
$V_n^{1/n}/N_n^{1/n}\to1$. The ratio of successive roots of $V$
along even $n$ would tend to $1/2$, contradicting eventual
nondecrease.
\end{proof}

Hence even a root-monotonicity substitute, up to vanishing multiplicative
error in the roots, is not necessary. The partition-product condition
is genuinely more flexible.

\subsection{Every finite distortion is realized}

\begin{theorem}[A one-parameter finite-loss family]
\label{thm:finiteloss}
For $c>1$, define
\begin{equation}
 N_0=1,\qquad N_n=n!\,c^{-n\boldsymbol1_{\{n\text{ odd}\}}}.
 \label{eq:finiteloss}
\end{equation}
Then $N_n^{1/n}\to\infty$, and its envelope is
\begin{equation}
 S_n=
 \begin{cases}
 n!,&n\text{ even},\\
 \max\{n!c^{-n},(n-1)!/c\},&n\text{ odd}.
 \end{cases}
 \label{eq:finitelossenvelope}
\end{equation}
In particular, $\Delta(N)=\mathcal I(N)=c$. The series $z-W_N$
has type one and inverse type exactly $c$.
\end{theorem}
\begin{proof}
The roots tend to infinity because $n!$ has superexponential roots and
the parity factor reduces them by at most $c$. One elementary lower
bound is $n!\ge\lfloor n/2\rfloor^{\lceil n/2\rceil}$.

For every composition of $n$, the product of the part factorials is at
most $n!$, and all powers of $c^{-1}$ reduce it. If $n$ is even, the
one-part composition attains $n!$. If $n$ is odd, a nontrivial
composition contains an odd part and hence a penalty of at least $c^{-1}$.
Also every nontrivial composition has factorial product at most
$(n-1)!$. To see this, group parts to two positive integers $a,n-a$;
then $a!(n-a)!=n!/\binom na\le(n-1)!$. The product before grouping is
no larger. Therefore every nontrivial product is at most $(n-1)!/c$,
attained by the partition $(n-1,1)$ for odd $n\ge3$. The one-part
product is $n!c^{-n}$. At $n=1$ the displayed maximum has the same
value $1/c$, so the formula remains valid.

For even $n$, $S_n/N_n=1$. For odd $n$,
\[
 \frac{S_n}{N_n}=\max\left(1,\frac{c^{n-1}}{n}\right),
\]
whose $n$th root tends to $c$ along the odd indices. Apply
Theorem~\ref{thm:distortionintro} and
Corollary~\ref{cor:extremaltype}.
\end{proof}

A finite distortion larger than one does not necessarily destroy the
zero-type class: \eqref{eq:twosideddistortion} still preserves type zero.
It does destroy exact positive type. The next example demonstrates a
stronger failure unavailable when the distortion is finite.

\subsection{Deep dips and an exact supermultiplicative repair}

\begin{theorem}[An infinite-distortion weight]
\label{thm:deepdips}
Let
\begin{equation}
 N_0=1,\qquad
 N_n=
 \begin{cases}
 2^{n^2},&n\text{ even},\\
 2^{(n^2-1)/4},&n\text{ odd}.
 \end{cases}
 \label{eq:deepdips}
\end{equation}
Then $N_n^{1/n}\to\infty$, but its envelope is
\begin{equation}
 S_n=
 \begin{cases}
 2^{n^2},&n\text{ even},\\
 2^{(n-1)^2},&n\text{ odd},
 \end{cases}
 \label{eq:deepdipsenvelope}
\end{equation}
and $\Delta(N)=\infty$.
\end{theorem}
\begin{proof}
The root exponents are $n$ on even indices and $(n^2-1)/(4n)$ on odd
indices, so both tend to infinity. An odd part $j\ge3$ in a partition
may be replaced by the even part $j-1$ and the part one without decreasing
the product, because $N_1=1$ and
\[
 (j-1)^2-\frac{j^2-1}{4}=\frac{(3j-5)(j-1)}4\ge0.
\]
Parts equal to one are left alone. After all such replacements, combine
the even parts into one even part: their squared exponents have sum at
most the square of their sum. If the original total $n$ is even, the
resulting product is at most $2^{n^2}$, attained by the one-part
partition. If $n$ is odd, at least one part one remains, so the product
is at most $2^{(n-1)^2}$, attained by $(n-1,1)$, or by $(1)$ when
$n=1$.

At odd $n$,
\[
 \frac1n\log_2\frac{S_n}{N_n}
 =\frac34n-2+\frac{5}{4n}\longrightarrow\infty.
\]
Thus the envelope repair has an unbounded exponential defect.
\end{proof}

\begin{theorem}[Type zero can invert to infinite type]
\label{thm:zerotoinfinity}
For the weight \eqref{eq:deepdips}, define
\begin{equation}
 f(z)=z-z^2-
 \sum_{\substack{n\ge2\\n\text{ even}}}
          2^{n^2-n\lceil\sqrt n\rceil}z^{n+1}.
 \label{eq:zerotoinfinityseries}
\end{equation}
Then $\Type(f)=0$ and $\Type(\inv f)=\infty$.
\end{theorem}
\begin{proof}
At even $n$, the normalized $n$th root of the coefficient of $f$ is
$2^{-\lceil\sqrt n\rceil}\to0$; at odd $n\ge3$ it is zero. The
single quadratic coefficient does not affect type. Hence $\Type(f)=0$.

Write the negative nonlinear coefficients of $f$ as $-b_{n+1}$ with
$b_{n+1}\ge0$. Its inverse has nonnegative coefficients by
Lemma~\ref{lem:lagrange}. For odd $n\ge3$, the two ordered compositions
$(n-1,1)$ and $(1,n-1)$ in the $k=2$ term yield
\begin{equation}
 [z^{n+1}]\inv f\ge
 (n+2)\,2^{(n-1)^2-(n-1)\lceil\sqrt{n-1}\rceil}.
 \label{eq:zerotoinfinitylower}
\end{equation}
Indeed the Lagrange factor is
$\binom{n+2}{2}/(n+1)=(n+2)/2$, and there are two distinct orderings;
also $b_2=1$. After division by $N_n$, the base-two logarithm of the
normalized $n$th root is bounded below by
\[
 \frac{\log_2(n+2)+(n-1)^2
 -(n-1)\lceil\sqrt{n-1}\rceil-(n^2-1)/4}{n}
 =\frac34n-\sqrt n+O(1),
\]
which tends to infinity along odd $n$. This proves the assertion.
\end{proof}

The mechanism is explicit: an even degree with a large weight mixes with
a single quadratic term and lands in a much smaller odd-degree weight.
The failure is not an oscillatory sign phenomenon; all inverse
contributions used in the lower bound are positive.
% ed. (2026-09-29): consequence for the polynomial test of the predecessor's thm:iff, left implicit by the author
\begin{ednote}
Since $N_n^{1/n}\to\infty$ for the deep-dip weight, the inverse of every
tangent-to-identity polynomial, whose coefficients grow at most
exponentially, has type zero; Theorem~\ref{thm:zerotoinfinity}
nevertheless gives a type-zero series whose inverse has infinite type.
Hence the equivalence of conditions (ii) and (iii) in Theorem~\texttt{thm:iff}
of \cite{repo-weight} (``the inverse of every tangent-to-identity polynomial
has type zero'') genuinely needs its log-convexity hypothesis. This is
consistent with that theorem, which assumes log-convexity.
\end{ednote}

\subsection{Harmless oscillations can be very irregular locally}

As a simple contrast, take
\begin{equation}
 N_n=(n+1)!,\qquad
 V_n=(n+1)!\,2^{(-1)^n\lfloor\sqrt n\rfloor}
 \quad(n\ge1),
 \label{eq:harmlessperturbation}
\end{equation}
with zeroth terms one. Proposition~\ref{prop:logconvex} makes $N$
admissible. Since $\log(V_n/N_n)=O(\sqrt n)=o(n)$,
Corollary~\ref{cor:perturbations} makes $V$ admissible with exactly the
same types. No direct check of adjacent quotients or log-convexity of
$V$ is required. Together with the preceding examples, this emphasizes
that the relevant scale is the envelope defect divided by $n$, not the
visible size of local oscillations alone.

\section{Consequences for coefficient transforms and transseries}
\label{sec:transseries}

\subsection{Exact radii of weight-normalized transforms}

For a scalar formal series with nonlinear coefficients $f_{n+1}$, define
its explicitly normalized transform
\begin{equation}
 \mathcal B_N f(\zeta)=
 \sum_{n\ge1}\frac{f_{n+1}}{N_n}\zeta^n.
 \label{eq:weighttransform}
\end{equation}
This definition is merely a coefficient normalization. It is not, for a
general $N$, a claim that an associated Laplace kernel or a summation
procedure exists.

The radius of convergence $R_N(f)$ is
\begin{equation}
 R_N(f)=\frac1{\Type(f)},
 \label{eq:radius}
\end{equation}
with $1/0=\infty$ and $1/\infty=0$. One can verify the formula directly:
inside the reciprocal limsup radius the terms are eventually bounded by
a convergent geometric series; outside it, a subsequence of terms fails
to tend to zero.

\begin{corollary}[Radius preservation and optimal radius distortion]
\label{cor:radii}
For admissible $N$, every tangent formal series satisfies
\begin{equation}
 R_N(\inv f)=R_N(f).
 \label{eq:radiusinvariance}
\end{equation}
If $N_n^{1/n}\to\infty$ and $1\le\Delta(N)<\infty$, then for
$0<R_N(f)<\infty$,
\begin{equation}
 \frac{R_N(f)}{\Delta(N)}\le R_N(\inv f)
                              \le\Delta(N)R_N(f).
 \label{eq:radiusdistortion}
\end{equation}
The constant is sharp. Radii zero and infinity are individually preserved
when $\Delta(N)<\infty$.
\end{corollary}
\begin{proof}
Take reciprocals in Theorems~\ref{thm:classificationintro} and
\ref{thm:distortionintro}, with the endpoint cases handled by the
zero- and infinite-type statements. The negative extremal input has
radius one and inverse radius $1/\Delta(N)$.
\end{proof}

For log-convex factorial or refined factorial coefficient scales this
recovers the radius interpretation that motivates the predecessor
\cite{repo-weight}. The extension is that the same interpretation now
has an exact necessary-and-sufficient criterion for arbitrary positive
scales, including the non-convex examples above.

\subsection{Where this fits in a transseries calculation}

An integer-indexed sector of a transseries calculation can produce a
formal series in a small parameter $q$, with a chosen coefficient scale.
When its normalization is tangent to the identity, the results here apply
to the reversion in $q$ without requiring the series to converge. They
also apply to finite systems of such formal series through
Section~\ref{sec:banach}. The envelope is an explicit preprocessing step:
compute or estimate $S_n$, and distinguish whether its defect is
sublinear, linear with finite limsup, or unbounded on the linear scale.

For example, a coefficient calculation known to have finite exact type
in the staircase scale \eqref{eq:staircase} can be inverted with precisely
the same type even though no equivalent log-convex scale exists. In the
family \eqref{eq:finiteloss}, a generic membership statement obscures a
real factor $c$ in the worst possible transformed radius. The deep-dip
example warns that a superexponential normalization by itself can hide
an arbitrarily severe nonlinear growth instability.
% ed. (2026-09-29): consequence for the near-linear feedback article's inversion question (editorial deduction, re-derived on filing)
\begin{ednote}
The repository article \path{Near_Linear_Boundary_Exponential_Feedback/}
asks for the exact root-growth scale of the compositional inverse of its
series $\widehat U=\sum u_nq^n$, for which $u_n^{1/n}\sim c_1\ee^{b_n}$
(its Theorem~\texttt{thm:main}); its editorial note says that the
log-convex theorem of \cite{repo-weight} would settle this if $\ee^{nb_n}$
were eventually log-convex. Log-convexity is not needed. That article's
Lemma~\texttt{lem:b} makes $b$ eventually nondecreasing, with
$b(x)\to\infty$ and $b(x)=o(\log x)$. Take the weight $M_m=\ee^{mb_m}$,
so $N_n=\ee^{(n+1)b_{n+1}}$, with $b_m$ given arbitrary positive values at
the finitely many small indices where it is undefined, and put
$V_n=\exp(n\max_{m\le n+1}b_m)$. Its roots are nondecreasing, so $V$ is
supermultiplicative, $V_0=1$, $V_n^{1/n}\to\infty$, and eventually
$\log(V_n/N_n)=-b_{n+1}=o(n)$. By
Theorem~\ref{thm:classificationintro}(iii) the weight is admissible;
$\Type(\widehat U)=c_1$, so Corollary~\ref{cor:scalarscaling} gives
$\Type(\inv{\widehat U})=1$, that is,
$\limsup_n|q_n|^{1/n}\ee^{-b_n}=1$ for the inverse coefficients $q_n$.
This is a limsup statement, not a root limit. The deduction is editorial;
it is not in either article.
\end{ednote}

These statements should not be extrapolated to an arbitrary Hahn support
by replacing the integer $n$ with an exponent. The proof uses finite
compositions of an integer and a concrete count of the possible lengths.
Accumulating action sets require a separate support-summability and
counting theory. Similarly, no Stokes constant, angular continuation,
Borel summability direction, or optimal truncation error follows merely
from the radius in \eqref{eq:radius}.

\section{Finite algorithms and executed verification}
\label{sec:verification}

\subsection{What a finite prefix can certify}

For rational input weights $N_1,\ldots,N_d$, the recurrence
\eqref{eq:envelopecalculus} computes $S_1,\ldots,S_d$ and records a
maximizing composition for each degree. It uses $O(d^2)$ rational
multiplications and comparisons. This is an arithmetic-operation count,
not a polynomial bit-complexity claim: the numerators and denominators
may be very large.

A second computation uses the ordered Lagrange formula, precomputing
truncated powers of $\sum_{j\ge1}N_jz^j$. An independent implementation
enumerates unordered partitions and evaluates
\eqref{eq:partitionformula}. Agreement of these paths checks the
coefficient index, multiplicity factor, and sign convention. A third
implementation computes a fixed number of formal Picard iterations and
checks the inverse equations. It does not use the partition formula.

For a chosen rational $\eta$ and $C$, an exact integer $B$ is found by
binary-searching the inequalities $B^j\ge N_j$ over the required finite
prefix. The quantities $P,Q,R$ of
Theorem~\ref{thm:finitecertificate} are then rational, so the certificate
can be tested with no transcendental arithmetic. Finite tests cannot
verify that $N_n^{1/n}\to\infty$ or that a limsup equals one; those
assertions are proved from the explicit hypotheses in the article.

\subsection{Recorded checks}

The accompanying standard-library Python program was executed with
Python 3.13.5, seed $20260929$, and maximum excess degree $36$ (ordinary
coefficient degree $37$). It reported \textbf{5,443 counted exact finite
checks}, all passing. The counter names and precise scope are preserved
in \texttt{data/verification.json}; a summary follows.

\begin{center}
\small
\begin{tabular}{@{}p{0.70\textwidth}r@{}}
\toprule
Check family & Count\\
\midrule
Envelope witnesses and positive inverse lower bounds & 216\\
Independent partition-formula comparisons & 108\\
Exact sparse--dense upper certificates & 648\\
Envelope supermultiplicativity inequalities & 3,780\\
Six weight families: two-sided inverse and Picard audits & 6\\
Explicit envelopes: two supermultiplicative families & 72\\
Deep-dip and finite-loss envelope formulas & 72\\
Additional arbitrary-weight finite audits & 144\\
Signed coefficient majorants and dilation checks & 288\\
Signed two-sided inverse cases & 24\\
Type-zero-to-infinite-type finite lower bounds & 17\\
Bivariate homogeneous coefficient-norm bounds & 48\\
Bivariate two-sided inverse cases & 8\\
Staircase convexity-defect identities & 12\\
\midrule
Total & 5,443\\
\bottomrule
\end{tabular}
\end{center}

The six structured families are the shifted factorial weight, the
staircase weight, the alternating supermultiplicative weight, finite loss
with $c=2$, deep dips, and the subexponential factorial perturbation.
Their positive inverse coefficients and envelopes were computed through
excess degree $36$. Independent partition enumeration was checked
through excess degree $18$. The two scalar inverse compositions were
checked through ordinary degree $16$ for these families, with independent
Picard checks through degree $12$. Twenty-four seeded signed-series
cases were checked through ordinary degree $13$.

Eight seeded bivariate maps were inverted through total degree $7$.
Both compositions were checked against the identity, and the homogeneous
coefficient-sum norms \eqref{eq:coefficientnorm} were compared with the
scalar inverse majorant. These are finite rational polynomial tests;
the infinite-dimensional statement follows from the multilinear proof,
not from sampling finite-dimensional cases.

Every asserted identity and inequality in the verification program uses
integers or rational numbers. The logarithms in the diagnostic tables
below use ordinary floating-point arithmetic, are labeled accordingly,
and are not interval certificates. A counted check may aggregate several
coefficient comparisons; the total is an audit counter, not a count of
independent mathematical theorems.

\subsection{A diagnostic separating envelope defect from tree overhead}

The finite observations are easiest to interpret after separating
\begin{equation}
 d_n=\frac1n\log_2(S_n/N_n),\qquad
 r_n=\frac1n\log_2(h_{n+1}(1)/S_n).
 \label{eq:diagnostics}
\end{equation}
The theorem predicts $r_n\to0$ for each of the displayed weights.
The envelope defect $d_n$ behaves very differently: it is zero for the
staircase, tends to one along odd indices in the $c=2$ finite-loss
family, and tends to infinity along odd indices for deep dips.

\begin{center}
\small
\begin{tabular}{@{}lrrr@{}}
\toprule
Weight & Excess degree $n$ & $d_n$ & $r_n$\\
\midrule
Staircase & 9 & 0.000000 & 0.154819\\
Staircase & 17 & 0.000000 & 0.066476\\
Staircase & 25 & 0.000000 & 0.138526\\
Staircase & 33 & 0.000000 & 0.032306\\
\addlinespace
Deep dips & 9 & 4.888889 & 0.384381\\
Deep dips & 17 & 10.823529 & 0.249878\\
Deep dips & 25 & 16.800000 & 0.190196\\
Deep dips & 33 & 22.787879 & 0.155433\\
\addlinespace
Finite loss, $c=2$ & 9 & 0.536675 & 0.507991\\
Finite loss, $c=2$ & 17 & 0.700737 & 0.273671\\
Finite loss, $c=2$ & 25 & 0.774246 & 0.200121\\
Finite loss, $c=2$ & 33 & 0.816837 & 0.161217\\
\bottomrule
\end{tabular}
\end{center}

These are rounded diagnostics from exact rational coefficients, not
certified logarithmic intervals. In particular, the staircase values of
$r_n$ are visibly nonmonotone. The proofs require no monotone convergence
and do not infer the limits from the displayed finite degrees.

\subsection{Reproduction and output discipline}

The package contains \texttt{article.tex}, its compiled PDF, the
verification code, recorded JSON and CSV data, provenance and proof-audit
notes, and build commands. The article embeds its bibliography and table;
it compiles without reading the verification outputs or accessing the
network. From the package directory, run
\begin{quote}\small
\begin{verbatim}
python3 code/verify.py --order 36
pdflatex -interaction=nonstopmode -halt-on-error article.tex
pdflatex -interaction=nonstopmode -halt-on-error article.tex
pdflatex -interaction=nonstopmode -halt-on-error article.tex
\end{verbatim}
\end{quote}
The verifier's default output is \texttt{build/verification/}, leaving
\texttt{data/} unchanged. An explicit \verb|--output| argument can
select another destination. It requires no third-party Python package,
contacts no external service, and does not modify the upstream repository.
PDF build and rendering checks are recorded separately from the
mathematical verification report.

\section{Proof dependencies and a formalization route}
\label{sec:formalization}

\subsection{The short dependency chain}

The universal classification rests on a particularly small chain of
arguments. First, finite Lagrange inversion converts inverse coefficients
to positive partition products. Second, replacing many small parts by
one large part proves that maximizing partitions have sublinear length.
Third, a sparse--dense decomposition bounds every product relative to
the supermultiplicative envelope and separates a small entropy cost from
an arbitrarily strong dense penalty. These three steps prove
\eqref{eq:central}. Dilation, absolute values, and the inverse involution
then prove the exact distortion and zero-loss criterion.

Composition introduces only the majorant
$B_C\circ B_C\preceq H_{2C}$. Unequal-type equality and all cancellation
examples use formal group identities. The Banach extension requires the
norm inequality for normalized symmetrized contractions, after which the
same scalar majorants apply. No compactness theorem for smooth functions,
analytic continuation theorem, or resurgent input is hidden in the chain.

\subsection{Finite algebra before asymptotics}

A practical formalization can begin with finite words of positive natural
numbers and their sum. The envelope recurrence, concatenation inequality,
and least-majorant property are finite statements about maxima. They can
be separated from formal power series entirely. The next layer is the
finite coefficient identity \eqref{eq:orderedformula}, together with
its absolute-value majorant and the equivalence to multiplicity-vector
sums.

The finite sparse--dense certificate is a useful standalone target. It
requires counting positive compositions and tracking the total mass of
parts at most $L$, but contains no limit. Formalizing it first isolates
the only combinatorial estimates needed for the asymptotic theorem.
The lower bound from a maximizing word should also be stated separately,
so that the $C<1$ case is not accidentally lost.

\subsection{Limits and extended-real endpoints}

The analytic layer can be divided into the sparse-maximizer lemma, the
binomial entropy estimate, the dense-suppression limit, and the central
root limit. One must preserve the order of quantifiers: first fix
$\eta$, then fix $B$ from finitely many weights, then let $n\to\infty$,
and only afterwards let $\eta\downarrow0$. The assertion is not uniform
in all $C>0$ simultaneously; each $C$ is fixed before the limit.

The type layer should make zero and infinity explicit. A convenient
approach is to prove finite coefficient bounds for every
$A>\Type(f)$, derive the upper inequality when the input type is finite,
and handle infinite types by applying the same inequality to the inverse.
The polynomial obstruction deserves its own lemma because it supplies
the necessity of a full root limit, not merely an unbounded subsequence.

\subsection{Vector-valued coefficients and what has not been checked}

For the Banach extension, a representation by symmetric continuous
multilinear coefficients avoids an implicit polarization loss. The
formalization needs normalized symmetrization, contraction, finite formal
composition, and the recursively defined inverse. Associativity can be
proved coefficientwise, independently of any convergent power-series
infrastructure.

No new Lean code is supplied or claimed to compile in this package.
The proposed stages are a dependency plan, not a statement that matching
library declarations already exist. The executed rational verification
program can provide regression examples after the formal statements are
implemented, but it does not certify the universal theorems.
% ed. (2026-09-29): existing repository Lean for the finite Lagrange layer
\begin{ednote}
The repository already contains a residue-free Lean proof of the finite
Lagrange--B\"urmann coefficient formula,
\path{Analysis/FabiusFunction/Lean/FabiusFunction/LagrangeInversion.lean}
(\texttt{Fabius.Lagrange.coeff\_solution}), and of the Catalan inverse of
$z+cz^2$ over any commutative ring,
\path{Analysis/FabiusFunction/Lean/FabiusFunction/QuadraticCompositionalInverse.lean}
(\texttt{QuadraticInverse.coeff\_succ\_inverse}); at $c=-1$ the latter is
the series used in Lemma~\ref{lem:polynomial}. Nothing about weights, types
or envelopes is formalized.
\end{ednote}

\section{Further research questions and concrete next targets}
\label{sec:questions}

The arbitrary-weight scalar zero-loss classification and the scalar
composition spectrum are settled by the proofs in this article. The
following questions concern additional information not encoded by those
answers. They are proposed research directions; not every item is
claimed to be a previously published open problem.

\subsection{Optimal subexponential overhead}

Theorem~\ref{thm:envelopeintro} proves
$\log h_{n+1}(C)-\log S_n=o(n)$, while
Theorem~\ref{thm:finitecertificate} supplies an explicit but generally
nonoptimal finite upper bound. Determine the smallest attainable
subexponential overhead from a specified weight profile. A concrete first
target is to express it in terms of the number and entropy of partitions
whose products lie within a chosen factor of $S_n$. This would refine
the repository's question on optimal reversion cost without presuming
log-convex local quotient asymptotics.
% ed. (2026-09-29): answered for factorial weights by a package of the same batch
\begin{ednote}
For factorial weights this is answered by the independent package
\path{Sharp_Subexponential_Cost_Gevrey_Reversion/} (filed in the same batch,
written without knowledge of this article). For
$N_j=((j+\nu)!/(1+\nu)!)^s$, $\nu\in\{0,1\}$, which is log-convex, so
$S=N$, its extremal series is $H_C$, and its main theorem gives
$\log h_{n+1}(C)-\log S_n\sim Cn^{1-s}$ for $0<s<1$, a multiplicative
equivalent for $s>1/2$ (with an extra constant at $s=1/2$), and a bounded
overhead for $s\ge1$, so the fixed coefficient ball is stable under
inversion exactly for $s\ge1$. The general-weight question stays open.
\end{ednote}

\subsection{Anisotropic exact types under a general Jacobian}

Theorem~\ref{thm:jacobian} proves that one scalar type and an invertible
matrix are insufficient for a unique inverse type. Define a directional
or multihomogeneous type profile that transforms exactly under linear
precomposition and is preserved, in an appropriate sense, by tangent
inversion. A useful finite-dimensional target is a profile indexed by
polyradii, together with an exact covariance law under diagonal
Jacobians. Non-diagonal changes would then test how much angular data
must be retained.

\subsection{Exact composition distortion outside the zero-loss class}

For a nonadmissible superexponential weight with finite $\Delta(N)$, the
same positive majorant gives an upper bound
$\Type(f\circ g)\le\Delta(N)\max(\Type(f),\Type(g))$. Is this
constant optimal? A single composition has a restricted two-level
substitution structure, whereas an inverse permits arbitrarily deep
structures. Identify the correct restricted envelope and determine
whether its defect can be strictly smaller than the inversion defect.
This asks for a new classification, not a repetition of the zero-loss
case.

\subsection{Effective criteria for presented weights}

The finite recurrence computes every prefix of $S$, but the condition
$\log(S_n/N_n)=o(n)$ is asymptotic. For which explicitly presented
families can it be decided or certified effectively? Candidate classes
include piecewise-polynomial logarithmic weights on residue classes,
regularly varying weights with controlled oscillation, and weights given
by finite max-plus recurrences. An algorithm should output a modulus of
the envelope defect, not extrapolate it from a long finite prefix.

\subsection{Quantitative cancellation profiles}

Theorem~\ref{thm:cancellation} determines everything that two scalar
types alone can say. What additional finite or asymptotic data predicts
the output type in a structured family of equal-type inputs? A concrete
target is a noncancellation theorem for prescribed sign cones or for
coefficients concentrated on separated degree sets. Another is a
stability theorem describing how a small coefficient perturbation moves
a pair away from the exact inverse relation.

\subsection{Differentiation, flows, and iterative logarithms}

Admissibility was characterized for inversion and composition only.
Determine the additional exact-type conditions needed for differentiation,
formal flows, and the iterative logarithm of a tangent map. Adjacent
weight ratios, absent from the main criterion, may reappear when an
operator changes degree. A basic test is the staircase weight: its exact
composition stability coexists with large local changes in its growth
profile. No implication about these other operations is assumed here.

\subsection{Hahn supports and accumulating actions}

Replace the integer excess degree by an ordered action semigroup. One
must specify which decompositions are finite, how their number is
controlled, and what notion of a coefficient-growth scale is invariant
under support refinements. A first target is a finitely generated positive
semigroup with a fixed positive additive degree. Only after that case is
understood should one address accumulating actions, where a finite
integer-composition count is no longer available.
% ed. (2026-09-29): repository packages on Hahn-support reversion
\begin{ednote}
For the support side of this question see the repository packages
\path{Support_Controlled_Reversion_One_Exponential/} (Hahn
Lagrange--B\"urmann inversion with a finite outer index bound at each
exponent), \path{Reversion_Beyond_Archimedean_Valuations/} (positive-word
support depth over an arbitrary ordered exponent group) and
\path{Action_Accumulation_Nonlinear_Inversion/} (inversion of
exponentially weighted action measures without a positive action gap).
None of them treats weighted coefficient types.
\end{ednote}

\subsection{Realization and summation compatible with the envelope}

For which admissible but non-log-convex weights do formal series admit
useful analytic realizations with bounds preserving the exact type?
Can a moment transform or acceleration procedure be built directly from
the envelope $S$, and when does its sum commute with inversion? These
questions require analytic continuation and growth hypotheses beyond
coefficient bounds. The equality of normalized-transform radii in
Corollary~\ref{cor:radii} supplies a necessary consistency check, not a
summation theorem.

\subsection{Finite certificates with bit-complexity bounds}

The present dynamic program has a quadratic count of rational operations,
while the integers in the examples can be huge. Develop compressed
representations for logarithmic or factorial weights and prove
bit-complexity bounds for envelope witnesses and sparse--dense
certificates. A useful target is an independently checkable certificate
whose size is small relative to the requested degree and the description
length of the weight, rather than to its expanded numerator.

\subsection{A proof-assistant formalization of the classification}

The finite maximum-product closure and the exact certificate offer a
manageable entry point. The main goal is a formally checked theorem
quantifying over all positive weights, including the necessity direction
and types zero and infinity. The Banach-space extension should then reuse
the scalar result through a single majorant interface. Formalizing just
selected numerical examples would not establish the classification;
formalizing the central root limit would be the decisive milestone.

\section*{Conclusion}
\addcontentsline{toc}{section}{Conclusion}

The obstacle to exact formal inversion is not the absence of
log-convexity. It is the exponential discrepancy between a weight and the
products that additive partitions can form from it. Superexponential root
growth suppresses the combinatorial cost of dense products, leaving the
supermultiplicative envelope as the exact exponential-scale answer.

This observation yields a complete arbitrary-weight criterion, an optimal
distortion constant, a canonical repair, and a sharp nonlinear type
calculus. The examples show that every part of the criterion matters:
strongly non-convex weights can work perfectly, finite parity penalties
can create a prescribed loss, and deeper parity dips can turn a
type-zero input into an inverse of infinite type. The same positive
majorants survive in formal Banach-space maps once the coefficient norm
is fixed precisely. What remains is not the basic scalar classification,
but finer overheads, directional information, different support geometries,
and analytic realization.

\begin{thebibliography}{9}

\bibitem{repo-weight}
ProveIt repository, maintained by Vladimir Reshetnikov.
\emph{Sharp Weighted Type Under Formal Reversion: Zero-loss inversion,
exact Borel radii, and non-D-finiteness in countable exponential feedback}.
Research article prepared 29 September 2026, including repository editorial
amendments of that date. Pinned file at commit
\nolinkurl{3d5973524506411392a911470b5ddc35521568ea}; Git blob
\nolinkurl{c2d9978806e0861f1bab7108b72ad92773e6a05c}.
Especially the research subsections ``Weights beyond log-convexity'',
``Two divergent composition arguments'', and ``Multivariate and
operator-valued inversion''.
\url{https://github.com/VladimirReshetnikov/ProveIt/blob/3d5973524506411392a911470b5ddc35521568ea/Analysis/Transseries/docs/series-and-transseries/Sharp_Weighted_Type_Formal_Reversion/article.tex}.
Accessed 29 September 2026.

\bibitem{gessel}
Ira M. Gessel.
\emph{Lagrange Inversion}.
arXiv:1609.05988, 2016.
\url{https://arxiv.org/abs/1609.05988}.
A survey of the classical coefficient formulas and their formal proofs.

\bibitem{rs-composition}
Armin Rainer and Gerhard Schindl.
\emph{Composition in ultradifferentiable classes}.
Studia Mathematica \textbf{224} (2014), no.~2, 97--131.
arXiv:1210.5102.
\url{https://arxiv.org/abs/1210.5102}.

\bibitem{rs-stability}
Armin Rainer and Gerhard Schindl.
\emph{Equivalence of stability properties for ultradifferentiable
function classes}.
Revista de la Real Academia de Ciencias Exactas, F\'isicas y Naturales.
Serie A. Matem\'aticas \textbf{110} (2016), no.~1, 17--32.
arXiv:1407.6673; the inspected preprint is version~2, 1 January 2015.
\url{https://arxiv.org/html/1407.6673v2}.

\end{thebibliography}
\end{document}
